%% Direction assignment from reaction chemistry (LLM ensemble).
%% \TBD values pending the full-database run, in progress.
\subsection{Direction assignment from reaction chemistry}\label{sec:methods-council}

Thermodynamic direction assignment is bounded by estimate coverage: a reaction
with no accepted \drG{} receives no call, regardless of how well characterised
its chemistry is. To assign direction for those reactions we evaluated a
complementary approach that reasons from the reaction itself rather than from
an energy estimate, using an ensemble of large language models.

\paragraph{Independence from the thermodynamic estimates.} Two prompt variants
were run. The heuristic-informed variant additionally presents the eQuilibrator,
group-contribution and dGPredictor estimates as evidence to weigh. The
heuristic-agnostic variant presents only the reaction identifier, its name and
its equation, with the equation rewritten to an undirected arrow so that no
direction is recoverable from the input text. Neither variant is shown the
database's own reversibility assignment. Only the heuristic-agnostic variant is
independent of the existing estimates, so only it can contribute a direction
call that is not a restatement of the numeric sources, and it alone is used
here. The informed variant is retained as a control: agreement between the two
bounds how much of the agnostic variant's output is recoverable from chemistry
alone.

\paragraph{Ensemble.} Each reaction is presented independently to three models
that do not share a provider, and each returns a direction
($>$, $<$, $=$, or $?$ where it judges the evidence insufficient), a confidence,
and a written justification grounded in the stoichiometry or in known enzyme
biology. A fourth model, not among the three, audits the resulting vote table
for internal contradiction. Where the three do not agree, a fifth model
adjudicates from the votes, the justifications and the audit; where they agree
unanimously and the auditor raises no objection, that answer is taken directly.
Model assignments were not chosen a priori. Eleven candidate models were
scored across the three roles against reactions carrying a measured
equilibrium constant, and the selected assignment is the one minimising the
rate at which a reaction is assigned the direction opposite to measurement,
which is the error class that propagates into flux predictions. The models
evaluated, their per-role scores and the selected assignment are given in
Supplementary \TBD[table ref].

\paragraph{Reaction population.} Of 56,012 reactions in the database we
exclude 7,609 flagged obsolete and 2,756 symmetrical transport
reactions, whose two sides are identical once compartment labels are removed
and which therefore carry no chemistry to reason about, leaving 45,647
reactions. Direction is assigned for all of them, including those the
thermodynamic pipeline already covers, so that the two methods can be compared
on the same population rather than only where the numeric sources are silent.

\paragraph{Validation and its limits.} Calls were compared against measured
equilibrium constants from the NIST Thermodynamics of Enzyme-Catalyzed
Reactions database, mapped to ModelSEED reaction identifiers through compound
aliases; 128 reactions in the evaluated slice carry such a measurement.
This is a small and reversibility-skewed sample, and it bounds what can be
claimed: it distinguishes clearly weaker configurations from stronger ones but
does not establish an accuracy figure with useful precision. Two limitations
follow and are reported rather than corrected. The ensemble commits to a
direction far more readily than the thermodynamic rules do, returning $?$ on
\TBD[pct]\% of reactions against the rules' much higher abstention rate, so its
errors are distributed differently: it assigns a direction where the rules
decline, and it occasionally assigns the opposite one. And its stated
confidence does not separate its correct calls from its incorrect ones, so
confidence cannot be used as an integration filter. Per-reaction justifications
and the auditor's objections are released alongside the calls
(Supplementary \TBD[data ref]), so that assignments can be reviewed
individually rather than accepted on aggregate.
